\JOUChapterWithEnglish{浮选柱气泡-颗粒作用}{Bubble-Particle Interaction in Column Flotation}
\label{chap:bubble-particle}

\JOUSectionWithEnglish{界面吸附与碰撞行为}{Interfacial Adsorption and Collision Behaviour}

微泡尺度、气液界面电性以及矿粒润湿特征共同决定了碰撞概率和黏附稳定性。相较于传统机械搅拌浮选，旋流-静态微泡浮选柱在低能耗条件下能够维持较高的气泡分散度，因此更适合处理细粒级和难浮颗粒体系。

\vspace{0.4\baselineskip}

通过界面吸附动力学分析可以发现，药剂浓度和停留时间的耦合作用会显著影响颗粒在界面上的再分布行为。该过程最终体现在精矿产率、灰分与循环负荷之间的协同变化上，为工艺优化提供了可解释的机理基础。
